\chapter{视频复原、增强与压缩}
\label{chap:vision-video-restoration}

夜间行车录像中的路牌在一帧里模糊，在下一帧里可能清晰一些。若能把同一处文字对应起来，邻帧就能补充当前帧缺少的观测；但错把车身纹理移到路牌上，也会产生比单帧更难察觉的错误。连续画面的处理因此既要利用时间冗余，也要控制运动、遮挡和传播带来的冲突。

\ref{chap:vision-image-restoration}章已经区分复原、呈现与压缩，\ref{chap:vision-video-recognition}章建立了帧间对应的基础。本章分别讨论如何取得有效跨帧信息、如何恢复或扩展时空采样，以及如何建立可解码的参考链。压缩可以直接处理输入视频，不要求先执行复原；如果二者组合，还需核验整个链路是否保持用途所需的内容。

\section{处理目标与观测条件的确定}
\label{sec:vision-video-restoration-goals}

“让视频更清楚”可能指去噪、改善亮度、增加分辨率或减少压缩损伤。不同目标允许改变的量不同，所需的输入和参照也不同。沿用图像复原章的记号，本章以$\bm{Y}_t$表示低质量观测，$\bm{X}_t$表示希望恢复的目标帧，$\widehat{\bm{X}}_t$表示算法输出；下标$t$是视频时间索引。

\subsection{退化观测与目标输出}
\label{sec:vision-video-degradation}

单帧中的噪声、模糊、低照度、分辨率和编码损失仍然存在，但它们可能随时间变化。自动曝光使同一路牌亮度改变，运动使车身边缘移动，遮挡使一块背景暂时不可见，切镜则可能让整个场景更换。只有在任务约定它们应被修正时，这些变化才是退化；正常车辆运动不应被算法抹平。

可以写$\bm{Y}_t=\mathcal{D}_t(\bm{X}_t)+\bm{\varepsilon}_t$概括该帧的成像与处理，其中$\mathcal{D}_t$包括时变模糊、采样或编码等作用，$\bm{\varepsilon}_t$表示额外噪声。这只是任务层面的简化，长曝光本身可能积分一段时间内的运动，不能总当作瞬时清晰图像上的固定卷积。多帧方法还需建立$\bm{X}_t$与其他时刻内容的对应，而不是默认所有目标帧都相同。

夜间录像中，平坦区域的随机颗粒提示噪声，移动路灯的拖尾可能来自曝光期间运动，整帧明暗变化可能来自曝光控制，方块边界则可能来自编码。针对不同成因，需要不同退化配对或数据标定。仅用一种合成高斯噪声训练的模型，不能据此宣称同时恢复所有这些变化。

\subsection{可用帧与等待范围}
\label{sec:vision-video-restoration-availability}

逐帧处理只利用当前帧，局部多帧处理读取邻域窗口，循环传播则把已经处理的信息压缩进状态。双向传播可以分别从前向后、从后向前累积信息，但后向分支必须先拿到未来帧。循环结构描述计算怎样复用，因果性描述输入何时可见，两者没有等价关系。

设输入帧率为$f$，重建当前帧要等待右侧$r$帧，则仅观测等待就至少约为$r/f$秒。每秒30帧、等待两帧时约为66.7毫秒，还未计入读取、推理和显示。离线方法若使用整段后向传播，应记录整段长度与缓存成本；把很长视频切成有限片段处理，又可能改变片段边界的恢复效果。

\subsection{内容、呈现与成本约束}
\label{sec:vision-video-content-cost}

观看用途可以允许一定色调变化，车牌读取要求字符保持，速度测量还要求时间戳和几何关系可追溯。模型生成一个看起来清晰的字符，并不说明该字符存在于观测中。扩展空间采样得到的高频细节、扩展时间采样得到的中间时刻，都需要与原始观测区分。

任务应同时约定分辨率、帧率、码率、等待、计算和存储预算。保存原视频、处理结果及所用帧范围，才能检查一次改变来自输入证据还是估计先验。若用户要求把红车改成蓝车或改变行动路线，问题已转为内容创作，见\ref{chap:vision-multimodal-generation}章。

\section{跨帧有效信息的取得}
\label{sec:vision-video-information-fusion}

邻帧的价值不是“多一张图”本身，而是它是否观察到目标位置需要的内容。取得这些信息通常要选择帧、确定对应位置、判断可见性，再聚合证据。位置正确与内容可靠是两个条件，曝光差异、反光和遮挡会使其中一个成立而另一个失效。

\subsection{邻帧与传播范围的选择}
\label{sec:vision-video-neighbor-selection}

固定窗口可以直接读取当前帧前后几张图，计算范围明确，却可能错过更早的清晰观测。单向传播让历史信息沿时间累积，同一模块复用于不同帧，但状态会压缩细节并传播错误。双向传播补充后侧信息，代价是等待未来及保存中间结果。

某块车牌在$t-3$帧清晰、$t$帧被雨滴遮住、$t+2$帧再次出现，局部窗口半径为1时就无法直接看到两次清晰观测。长时传播可能把较早特征带来，但不能保证车牌细笔画仍保存在状态中。若中间切镜，继续传播旧场景特征反而有害，可按具体系统的切镜规则重置状态。扩大范围后的收益要与这些条件共同检查。

\subsection{对齐与可见性的判断}
\label{sec:vision-video-warping}

设目标帧位置为$\bm{p}=(x,y)^\top$，由目标帧$t$指向邻帧$s$的光流为$\bm{u}_{t\to s}(\bm{p})$。为了在目标网格中重建同一内容，应到邻帧的位置$\bm{p}+\bm{u}_{t\to s}(\bm{p})$取值。这个方向约定很重要：若误用$s\to t$的位移而不转换，通常不能靠简单取负得到正确采样位置。

\term{反向取样}（backward warping）逐个遍历输出位置，到源图像查询内容，避免向前移动源像素时形成空洞或多对一写入冲突。查询位置通常不是整数。若横纵坐标分别为$j+\alpha,i+\beta$，其中$0\leq\alpha,\beta\leq1$，双线性插值用周围四个像素计算某个通道的值：
\begin{equation}
  \begin{aligned}
    \widetilde y={}&(1-\alpha)(1-\beta)y_{i,j}
       +\alpha(1-\beta)y_{i,j+1}\\
       &+(1-\alpha)\beta y_{i+1,j}
       +\alpha\beta y_{i+1,j+1}\text{。}
  \end{aligned}
  \label{eq:vision-video-bilinear}
\end{equation}
这里的四个$y$都来自邻帧$s$同一通道，省略了时间与通道标记。权重随距离连续变化且和为1，但这只能解释插值，不保证邻帧位置的内容确实可见。

在教学网格中，四角值从左上到右下为$0.2,0.6,0.4,1.0$，取样位置为$\alpha=0.25,\beta=0.5$。对应权重为$0.375,0.125,0.375,0.125$，得到
\[
  \widetilde y=0.375(0.2)+0.125(0.6)
       +0.375(0.4)+0.125(1.0)=0.425\text{。}
\]
图\ref{fig:vision-bilinear-warp}把位置与权重放在同一网格中。若移动字母的位移估计少了一个像素，取来的可能是字母边缘与背景的混合，和当前帧融合后就会产生双边缘。

\begin{figure}[htbp]
  \centering
  % 四角与查询点均为正文可复算的教学数值。
\begin{tikzpicture}[x=1cm,y=1cm,font=\figurefont\small]
  \draw[BookRule,line width=0.6pt] (1,0) rectangle (4.5,-2.8);
  \foreach \x/\y/\val/\wt/\anchor in {
    1/0/0.2/0.375/south west,
    4.5/0/0.6/0.125/south east,
    1/-2.8/0.4/0.375/north west,
    4.5/-2.8/1.0/0.125/north east}{
    \fill[BookBlue] (\x,\y) circle (2pt);
    \node[anchor=\anchor] at (\x,\y) {\val\ (\wt)};
  }
  \draw[BookMuted,dashed,line width=0.6pt] (1.875,0) -- (1.875,-2.8);
  \draw[BookMuted,dashed,line width=0.6pt] (1,-1.4) -- (4.5,-1.4);
  \fill[BookTeal] (1.875,-1.4) circle (2.5pt);
  \node[above right] at (1.875,-1.4) {取样点};
  \draw[-{Stealth[length=2mm]},BookInk,line width=0.8pt] (0.6,0.65) -- (4.8,0.65);
  \node[anchor=west] at (4.8,0.65) {$x$};
  \draw[-{Stealth[length=2mm]},BookInk,line width=0.8pt] (0.5,0.45) -- (0.5,-3.2);
  \node[below] at (0.5,-3.2) {$y$};
  \node[anchor=west,align=left] at (5.7,-0.6) {$\alpha=0.25$\\$\beta=0.5$};
  \node[anchor=west,align=left] at (5.7,-2.0) {四角加权\\$\widetilde y=0.425$};
\end{tikzpicture}

  \caption{双线性取样的教学网格。横坐标向右、纵坐标向下；查询点位于两行的中间、偏左四分之一处。四角旁括号内是插值权重，权重和为1，输出为0.425。图中线段表示插值邻域，不表示物体真实运动轨迹。}
  \label{fig:vision-bilinear-warp}
\end{figure}

图像对齐移动像素值，特征对齐移动网络提取的表示。两者都要统一坐标和尺度；低分辨率特征上的一个单位，不一定等于原图一个像素。越界位置需有明确的填充及有效性标记，遮挡位置则不能仅凭插值有数值就参与可靠融合。前后向光流一致性可以提供筛选线索，但自身也会受流估计误差影响。

即使坐标对应正确，不同曝光和反光也可能改变亮度。可见性、光度校正或可靠性权重因此属于额外判断。多位置取样允许在初始对应附近寻找信息，但要说明偏移由谁预测、权重怎样形成，不能把“对齐模块”当作已经完成这两项计算。

\subsection{有效信息的聚合与传播}
\label{sec:vision-video-fusion}

考虑当前观测值0.40与对齐后的邻帧值0.425，假设可靠性分别为1和3，加权合成为$(0.40+3\times0.425)/4=0.41875$。若错误对应带来0.90且仍赋予权重3，结果就变成0.775。这个教学对照说明，多帧汇总是否有益取决于对应及可靠性，分母归一化并不能消除错配。

实际网络通常聚合特征而不直接平均像素，并把聚合结果继续传播。当前特征既要保留本帧观测，也要有机会修正历史错误。下面的RVRT以短片段为单位联合更新，\ref{sec:vision-video-basicvsr}节再展开逐帧双向传播，两者共享问题但组织方式不同。

\subsubsection{RVRT的局部片段处理与引导多位置聚合}
\label{sec:vision-video-rvrt}

RVRT将低质量视频划分为固定长度的相邻片段，在片段内让多帧特征通过时空窗口注意力共同更新，再把前一个片段的状态对齐后传给后一个片段。这样既保留片段内部的联合处理，又通过循环状态连接更长时间\scite{vision-rvrt2022}

跨片段对齐需要为当前某个位置，在支持片段的多帧中寻找相关内容。\term{引导可变形注意力}（guided deformable attention, GDA）先用光流提供初始取样位置，再根据当前特征、预对齐支持特征及光流预测多个偏移。在支持帧$s$中，第$m$个位置可写为
\[
  \bm{p}_{s,m}=\bm{p}+\bm{u}_{t\to s}(\bm{p})
                   +\Delta\bm{p}_{s,m}\text{。}
\]
每个支持帧提供多个位置，非整数处沿用双线性取样。光流缩小搜索范围，偏移调整在哪取值，注意力再决定每个值占多大比重。

令当前查询为$\bm{q}$、第$(s,m)$个采样键和值为$\bm{k}_{s,m}$、$\bm{v}_{s,m}$，单个注意力头计算
\begin{equation}
  \begin{aligned}
    a_{s,m}&=
      \frac{\exp(\langle\bm{q},\bm{k}_{s,m}\rangle/\sqrt d)}
      {\sum_{s',m'}\exp(\langle\bm{q},\bm{k}_{s',m'}\rangle/\sqrt d)},\\
    \bm{r}&=\sum_{s,m}a_{s,m}\bm{v}_{s,m}\text{，}
  \end{aligned}
  \label{eq:vision-video-gda}
\end{equation}
其中$d$是查询和键的通道维度。RVRT用前一细化层的特征形成查询与键，用当前层已传播的支持特征形成值，使比较依据与待传递内容各有来源。多头或分组分别执行取样与聚合，再组合通道信息；新的片段状态还结合该片段之前各层的特征共同细化。

车牌边缘的光流初始位置可能偏向笔画一侧。多位置取样可以同时覆盖两侧及其他支持帧，再按相似度聚合，但若所有采样点都落在雨滴上，注意力也没有真实字符可读。原方法还用预测偏移的平均值逐层修正光流引导，并在不同细化层交替传播方向。因此，网络有循环结构，仍会使用当前帧之后的信息。

训练以低质量序列和高质量目标配对，使用平滑绝对误差形式的Charbonnier重建损失学习取样、聚合及重建。原论文验证了超分、去模糊和去噪，这些结果不直接证明低照度增强效果。对切镜或错误传播，应分别查看帧内清晰度与后续连续性；重复传递的错误特征不会因为来自多个状态就变成多份独立证据。

\section{目标画面的恢复与采样扩展}
\label{sec:vision-video-restoration-output}

取得多帧信息后，输出可以保持原来的空间与时间网格，也可以增加空间像素或插入中间时刻。网格变化决定哪些位置有直接观测、哪些必须依靠估计。分辨率和帧率同时增加时，两种估计条件要一起满足。

\subsection{原采样位置上的恢复与增强}
\label{sec:vision-video-native-restoration}

同分辨率去噪或去模糊仍在原帧时刻输出，邻帧只提供补充证据。训练样本应说明目标是怎样取得的，例如对清晰序列合成规定噪声，或利用已对齐的模糊与清晰拍摄配对。目标帧错位会让网络受到互相冲突的像素监督，合成退化与真实夜间成像不一致也会影响迁移。

RVRT的去噪、去模糊配置沿用前节的跨片段机制，但在浅层先降低特征分辨率以控制计算，再重建回原空间大小；输出尺度为1，不意味着所有中间特征都保持原尺寸。是否需要改变输入通道、噪声条件或训练数据，要对应具体配置，不能只凭网络名称判断用途。

夜间路牌提亮后，检查的不仅是平均亮度，还包括相邻帧是否反复明暗跳变、白色字符是否偏色、运动边缘是否拖尾。降噪抹掉细笔画、锐化改变字符、局部增强放大闪烁，都是同采样输出仍可能失真的例子。增强服务观看时可以接受的改变，用于文字读取或亮度测量时可能不可接受。

\subsection{空间采样扩展与细节恢复}
\label{sec:vision-video-super-resolution}

相邻帧若有亚像素位移，可能观察到不同的采样位置。为理解这一点，考虑一个简化的一维清晰信号$\bm{x}=(x_1,x_2,x_3,x_4)^\top$。第一帧只采到$x_1,x_3$，第二帧因半个低分辨率像素的位移采到$x_2,x_4$。在位置精确、没有模糊和噪声的教学条件下，两帧合起来给出四个分量，而单帧仅给出两个。

真实相机通常还进行光学模糊和像素积分，观测不是任意选出几个高分辨率点。更完整的关系包含模糊、位移及下采样，即使多帧合起来也未必确定所有高频细节。遮挡、流估计偏差和颜色变化还会破坏互补条件。多帧超分可以利用观测差异，不能据此宣称所有新增像素都已经被真实拍摄。

\subsubsection{BasicVSR的双向传播、特征对齐与上采样}
\label{sec:vision-video-basicvsr}

BasicVSR以相邻低分辨率帧的光流对齐特征，分别沿前向和后向传播状态，同一时刻再把两路信息拼接用于上采样。原方法把传播、对齐、聚合和上采样组织成清晰的计算链，便于逐项分析视频超分的信息来源\scite{vision-basicvsr2021}

以前向分支为例，估计从当前帧$t$指向前帧$t-1$的光流，把前一状态取样到当前坐标，再与当前低清帧拼接更新。后向分支对称地读取$t+1$。用$\mathcal{W}$表示取样、$R_{\to}$和$R_{\leftarrow}$表示两路更新网络：
\begin{equation}
  \begin{aligned}
    \bm{h}^{\to}_t&=
      R_{\to}\!\left(\bm{Y}_t,\,
      \mathcal{W}(\bm{h}^{\to}_{t-1},\bm{u}_{t\to t-1})\right),\\
    \bm{h}^{\leftarrow}_t&=
      R_{\leftarrow}\!\left(\bm{Y}_t,\,
      \mathcal{W}(\bm{h}^{\leftarrow}_{t+1},\bm{u}_{t\to t+1})\right)\text{。}
  \end{aligned}
  \label{eq:vision-video-basicvsr}
\end{equation}
序列两端缺少的状态初始化为零。原实现用光流网络估计相邻输入的对应，用卷积和残差块更新状态；两路在各自方向上传播，不把后续BasicVSR++的二阶传播和对齐设计混入这一版本。

同一时刻把$\bm{h}^{\to}_t$与$\bm{h}^{\leftarrow}_t$沿通道拼接，再经卷积和像素重排形成高分辨率结果。像素重排把$h\times w\times(r^2c)$的特征重新组织为$rh\times rw\times c$，其中$r$为放大倍数。它改变数据排列；新增高分辨率值由前面的学习计算产生，而非重排这一动作从无到有地恢复细节。实现还可加回输入的插值基底以学习高分辨率残差。

训练使用低清序列与高清目标，通过逐元素Charbonnier损失约束重建。对误差$e$，平滑惩罚为$\rho_\epsilon(e)=\sqrt{e^2+\epsilon^2}$，再对时空位置与通道取平均；$\epsilon>0$控制零附近的平滑程度。训练时改变编码、对齐和重建参数，使用时只运行这些模块与状态传播，不针对每段视频重新拟合网络。

三帧车牌中，第一帧清楚的左侧笔画通过前向状态传到第二帧，第三帧清楚的右侧笔画通过后向状态补充第二帧。去掉对齐，同一位置可能来自不同笔画；去掉传播，只剩局部或单帧证据；去掉上采样，输出网格不会增加。拼接只让两路信息同时可用，是否利用正确仍由训练和输入条件决定。

图\ref{fig:vision-basicvsr-propagation}把两条时间依赖与一个前向更新分别画出，可以沿取样、拼接和更新检查历史怎样进入当前状态。
\begin{figure}[htbp]
  \centering
  % 上半图省略每列共有的输入/重建边，以突出两条时间依赖。
\begin{tikzpicture}[x=1cm,y=1cm,font=\figurefont\small,
  state/.style={draw=BookTeal,fill=BookTeal!10,minimum width=17mm,
    minimum height=7mm,line width=0.8pt},
  block/.style={draw=BookInk,fill=BookPaper,minimum width=18mm,
    minimum height=8mm,align=center,line width=0.8pt},
  arr/.style={-{Stealth[length=2mm]},BookInk,line width=0.8pt}]
  \node[anchor=west] at (0,4.8) {两条独立的时间传播};
  \foreach \x/\i in {1.3/1,5/2,8.7/3}{
    \node at (\x,4.2) {时刻$\i$};
    \node[state] (f\i) at (\x,3.5) {$\bm h^{\to}_{\i}$};
    \node[state] (b\i) at (\x,2.4) {$\bm h^{\leftarrow}_{\i}$};
  }
  \draw[arr] (f1) -- (f2);
  \draw[arr] (f2) -- (f3);
  \draw[arr] (b3) -- (b2);
  \draw[arr] (b2) -- (b1);
  \node[anchor=west] at (0,1.45) {一个前向更新};
  \node[block] (old) at (1.1,0.6) {历史状态\\$\bm h^\to_{t-1}$};
  \node[block] (warp) at (4.1,0.6) {取样对齐};
  \node[block] (concat) at (7.0,0.6) {通道拼接};
  \node[state] (new) at (9.4,0.6) {$\bm h^\to_t$};
  \draw[arr] (old) -- (warp);
  \draw[arr] (warp) -- (concat);
  \draw[arr] (concat) -- node[above] {$R_\to$} (new);
  \node (flow) at (4.1,-0.6) {$\bm u_{t\to t-1}$};
  \node (input) at (7.0,-0.6) {$\bm Y_t$};
  \draw[arr] (flow) -- (warp);
  \draw[arr] (input) -- (concat);
\end{tikzpicture}

  \caption{BasicVSR的时间依赖与单步更新。上半图只画两路状态传播，每列都使用该时刻输入帧，同一列两路状态共同用于重建。下半图展开一个前向更新：历史状态按光流取样，再与当前帧拼接并更新。后向分支必须等未来帧可用。}
  \label{fig:vision-basicvsr-propagation}
\end{figure}
\FloatBarrier

双向方法依赖未来。若改成单向在线运行，结构、训练及边界条件都发生变化，应重新验证，不能只删除后向结果就沿用原成绩。大运动、遮挡和切镜可能使历史特征拖入错误区域，即使车牌看上去变清晰，也须回查字符是否有观测支持。

\subsection{时间采样扩展与中间帧估计}
\label{sec:vision-video-interpolation}

\term{视频插帧}（video frame interpolation）根据相邻或相隔一定时间的观测，估计两者之间的图像。它有两侧输入；只依据历史预测未来缺少后一侧约束，属于不同设置。缺帧补全也要注明缺失长度和实际可用邻帧，不能假设缺一帧与缺半秒具有相同难度。

\subsubsection{FILM的多尺度运动估计与中间帧合成}
\label{sec:vision-video-film}

两辆车在输入帧之间移动很多像素时，单一尺度的局部匹配可能找不到同一边缘。FILM在图像金字塔上提取具有可比含义的多尺度特征，从粗到细估计中间时刻到两侧输入的取样位移，再将对齐的图像与特征共同交给合成网络。

特征提取器在相同深度的跨尺度处理间共享权重，并把空间分辨率相同的不同深度特征组合起来，使运动估计可以复用跨尺度信息。粗尺度上40像素运动缩成较小位移，较容易进入局部比较范围；把粗位移上采样后，细尺度网络再预测增量。位移若以当前网格像素为单位，从半分辨率升到全分辨率时，除了扩展网格，还要把位移数值乘2。

FILM直接估计中间帧到两侧的任务导向位移，而不先要求给出真实中间图像。模型只看两端图像，依据其特征预测在哪里取值；训练通过最终中间帧监督约束这一路径。因此，不需要额外的光流或深度真值，但预测位移也不能据此当作独立验证的真实光流\scite{vision-film2022}

若中间时刻为$\tau=1/2$，一辆车在两端横坐标20和28之间匀速平移，教学中间位置为24。输出位置24应分别向第一帧的20和第二帧的28取样，位移为$-4$和$+4$。这说明两路取样方向。真实模型还要处理加速、遮挡及新显露背景，不是对两个输入坐标机械取平均；两侧都没见过的区域只能由合成先验估计。

各尺度上对齐后的图像、特征和位移共同进入类似U-Net的解码器，输出中间帧。训练样本是两端与中间真实帧组成的三元组。原工作区分仅用像素绝对误差的配置，以及组合像素、VGG特征与Gram统计损失的配置；后者的特征与统计约束见\ref{sec:vision-generation-perceptual-style}节。不同目标可能改变锐利程度与像素误差，报告结果时应注明使用哪一版本。

原论文主要以中点三元组训练和评价。需要更密时间网格时，可反复插入中点，得到$1/4,1/2,3/4$等位置；新一轮可能把先前生成帧作为输入，误差也会继承，不能当作新增实拍约束。任意时刻的直接插值若依赖额外时间条件或不同实现，需要另外说明。

在线插帧必须等待后一端。30帧每秒输入若估计两个原帧之间的中点，中点内容至少在后一帧到达后才可能生成；相对该中点的观测等待约为16.7毫秒，若维持整体缓冲播放则还要记录采用的缓冲延迟和计算耗时。平滑慢动作适合呈现，却不能证明插出的碰撞、接触或遮挡细节真实发生过。

\section{有限比特下的视频编码与解码}
\label{sec:vision-video-codec}

相邻帧通常有大量相同内容，传输每帧完整表示会反复编码这些信息。帧间编码利用已经解码的图像预测当前帧，再编码预测尚未解释的部分。这里把待编码的输入仍记为$\bm{Y}_t$，解码重建记为$\widehat{\bm{Y}}_t$；压缩目标是保持这个实际输入，不默认输入已是前节所追求的无退化真值。

\subsection{参考依赖与待编码内容的确定}
\label{sec:vision-video-codec-reference}

帧内编码独立描述一帧，可用作首帧或重建参考的起点；帧间编码依赖其他已重建帧。编码器能看原视频，解码器只能读码流和此前重建的参考。因此，计算帧间预测时，两端必须使用相同的已解码参考，不能让编码器偷偷使用解码端尚不可得的原图。

\subsubsection{DVC的运动补偿与残差参考链}
\label{sec:vision-video-dvc}

DVC把运动估计、运动压缩、补偿预测和残差压缩组织为可联合训练的网络。它以当前源帧$\bm{Y}_t$和上一重建帧$\widehat{\bm{Y}}_{t-1}$估计光流，再把光流编码为潜在表示、量化并解码，取得实际可传输后恢复的位移$\widehat{\bm{u}}_t$\scite{vision-dvc2019}

运动补偿使用解码后的位移对参考帧取样，并把取样结果、参考帧和位移一起交给补偿网络，生成预测帧$\overline{\bm{Y}}_t$。原始估计光流若未经编码量化，解码端无法取得它；用它产生训练或测试预测会掩盖运动压缩误差。补偿预测完成后，才计算当前源帧还差什么：
\begin{equation}
  \bm{E}_t=\bm{Y}_t-\overline{\bm{Y}}_t,\qquad
  \widehat{\bm{Y}}_t=\overline{\bm{Y}}_t+\widehat{\bm{E}}_t\text{。}
  \label{eq:vision-video-dvc-reconstruction}
\end{equation}
$\bm{E}_t$是编码端可算的残差，$\widehat{\bm{E}}_t$是残差经编码、量化与解码后的重建。当前重建帧写入参考缓存，成为下一帧的共同条件。图\ref{fig:vision-dvc-reference}按两端职责画出这条依赖。

\begin{figure*}[htbp]
  \centering
  % 原图仅编码端可见；下路也由编码端本地执行，保证共同参考。
\begin{tikzpicture}[x=1cm,y=1cm,font=\figurefont\small,
  block/.style={draw=BookInk,fill=BookPaper,minimum width=24mm,
    minimum height=9mm,align=center,line width=0.8pt},
  arr/.style={-{Stealth[length=2mm]},BookInk,line width=0.8pt,rounded corners=1mm}]
  \node[anchor=west] at (-0.2,4.7) {编码端：从原帧与本地预测形成两类符号};
  \node[block] (srcm) at (1.4,3.6) {$\bm Y_t,\widehat{\bm Y}_{t-1}$};
  \node[block] (encm) at (5.2,3.6) {运动估计\\编码、量化};
  \node[block] (symm) at (9.0,3.6) {$\widehat{\bm m}_t$};
  \node[block] (srcr) at (1.4,2.2) {$\bm E_t=\bm Y_t-\overline{\bm Y}_t$};
  \node[block] (encr) at (5.2,2.2) {残差编码\\量化};
  \node[block] (symr) at (9.0,2.2) {$\widehat{\bm z}_t$};
  \draw[arr] (srcm) -- (encm);
  \draw[arr] (encm) -- (symm);
  \draw[arr] (srcr) -- (encr);
  \draw[arr] (encr) -- (symr);
  \draw[BookTeal,fill=BookTeal!8,line width=0.8pt] (12.7,1.65) rectangle (15.6,4.15);
  \node at (14.15,3.6) {运动码流};
  \node at (14.15,2.2) {残差码流};
  \draw[arr] (symm.east) -- node[above] {熵编码} (12.7,3.6);
  \draw[arr] (symr.east) -- node[above] {熵编码} (12.7,2.2);
  \draw[BookRule,line width=0.6pt] (-0.2,1.2) -- (15.6,1.2);
  \node[anchor=west] at (-0.2,0.6) {解码端：先从码流恢复符号};
  \node[block] (dm) at (1.4,-0.2) {$\widehat{\bm m}_t$};
  \node[block] (decm) at (5.2,-0.2) {运动解码};
  \node[block] (comp) at (9.0,-0.2) {运动补偿};
  \node[circle,draw=BookTeal,fill=BookTeal!8,line width=0.8pt,
    minimum size=6mm,inner sep=0pt] (sum) at (12.2,-0.2) {$+$};
  \node[block] (out) at (14.4,-0.2) {$\widehat{\bm Y}_t$};
  \draw[arr] (dm) -- (decm);
  \draw[arr] (decm) -- node[above] {$\widehat{\bm u}_t$} (comp);
  \draw[arr] (comp) -- node[above] {$\overline{\bm Y}_t$} (sum);
  \draw[arr] (sum) -- (out);
  \node (ref) at (9.0,0.7) {$\widehat{\bm Y}_{t-1}$};
  \draw[arr] (ref) -- (comp);
  \node[block] (dr) at (1.4,-1.7) {$\widehat{\bm z}_t$};
  \node[block] (decr) at (5.2,-1.7) {残差解码};
  \draw[arr] (dr) -- (decr);
  \draw[arr] (decr.east) -- node[above,pos=0.4] {$\widehat{\bm E}_t$}
    (12.2,-1.7) -- (sum.south);
  \node[align=center] (buffer) at (14.4,-1.7) {保存为\\下一帧参考};
  \draw[arr] (out.south) -- (buffer.north);
\end{tikzpicture}

  \caption{DVC式单参考帧间编码的主要依赖，根据原论文的计算职责重绘。上路在编码端形成量化运动与残差符号并写入码流，下路说明独立解码器如何从码流重建。两端运动补偿均使用解码后的运动和上一重建帧，编码器还须执行与下路相同的本地解码以更新参考；原始当前帧只在编码端可见。}
  \label{fig:vision-dvc-reference}
\end{figure*}
\FloatBarrier

用四个像素表示横向移动方块。上一重建为$(0,1,1,0)$，向右移动一个像素、左边以零填充的预测为$(0,0,1,1)$。设当前源帧为$(0.2,0,1,1)$，左侧新进入背景并非零，运动预测无法解释它；残差为$(0.2,0,0,0)$。若残差解码得到$(0.18,0,0,0)$，当前重建就是$(0.18,0,1,1)$，四像素均方误差为$0.02^2/4=10^{-4}$。这是教学简化，省略了补偿网络的额外修正。

下一帧要使用含0.18的重建参考，而不能使用编码端掌握的原值0.2。否则，两端即使拿到相同运动和残差，也会从不同起点产生不同图像，误差继续沿链传播。首帧需要另行编码，DVC的帧间网络并没有消除这个起点成本。

\subsection{运动、残差与辅助信息的编码}
\label{sec:vision-video-rate-allocation}

运动编码器把稠密位移变成潜在表示$\bm{m}_t$，残差编码器把$\bm{E}_t$变成潜在表示$\bm{z}_t$。它们分别量化为$\widehat{\bm{m}}_t$和$\widehat{\bm{z}}_t$，再根据概率模型进行熵编码。接收端先把比特还原为相同符号，再运行运动和残差解码网络。潜在张量的元素数并不等于码流大小，实际长度还受概率分布和编码格式影响。

联合训练在失真与两类表示的码率之间作取舍。对一帧可用
\begin{equation}
  L_t=\lambda D(\bm{Y}_t,\widehat{\bm{Y}}_t)
       +R_{\mathrm{motion},t}+R_{\mathrm{residual},t}
  \label{eq:vision-video-rate-distortion}
\end{equation}
表示DVC的主要目标，其中$D$在原配置中取均方误差，$R$是由量化表示概率估计的编码成本，$\lambda>0$控制质量权重。若$D$按像素平均、$R$用每帧比特，则$\lambda$的尺度与改用比特/像素时不同，比较时应固定归一化。

训练中的量化采用加均匀噪声的可导近似，使用时则真正舍入为离散符号。概率模型以符号负对数概率估计比特需求，相关原理见\ref{sec:vision-restoration-bitstream}节。运动估计也参与率失真目标，因而其输出需要兼顾可压缩性与预测用途，不应仅以光流端点误差评价整个编码系统。

为处理训练时的参考依赖，原方法缓存重建帧供后续样本使用，随训练更新缓存，避免每次都从序列首帧回算整条链。这是训练组织方式。真正编码时仍要顺着当前参数形成的重建链运行，独立解码器不能依赖训练缓存或原视频。

假设同一帧的方案A花20比特描述运动、80比特描述残差，方案B花10比特描述运动、105比特描述残差，总数分别为100与115。运动分支省了10比特，整个系统却多花15比特，因为粗运动把更多误差留给残差。这个教学对照解释联合目标的必要性，并不假设两方案失真一定相同；若失真不同，还要一起报告。

完整视频应统计首帧或关键帧、各帧运动、残差以及实际传输的尺寸、模式和其他辅助信息。估计码率用于训练，真实码流字节数用于核验传输成本。模型若预先固定在两端，需说明模型版本是否计入部署成本；若传输中更新模型或概率参数，那些新增字节也属于必要信息。

\subsection{按依赖关系解码与重建}
\label{sec:vision-video-decoding}

独立解码器先恢复首帧，再按参考关系读取后续帧的运动与残差符号，重建位移和残差，执行补偿与相加，然后把当前结果写回参考缓存。所有运算的尺寸、裁剪、量化和数值约定都应与编码器的本地解码一致。只有从保存的码流独立重建成功，才能说明输出是一套可用的编码结果。

单参考链通常不能从任意中间帧直接解码，随机访问需要从相应关键帧开始恢复依赖。关键帧更密会增加独立帧的成本，却缩短随机访问等待及错误传播范围。传输丢包还可能破坏后续参考，不能只检查每帧单独的解码函数。多参考或双向编码可以改变依赖与压缩效率，也会改变缓冲、排序及延迟条件，应作为不同编码配置比较。

面向观看可检查纹理和闪烁，面向识别还需核验对象、车牌和路标。DVC是学习式单参考编码的代表，原论文与传统标准的数值比较对应特定实现和参数；编码标准名并不能替代具体编码器、预设、关键帧间隔和速率控制配置，更不能据历史实验给出当前持续有效的排名。

\section{处理效果的验证与取舍}
\label{sec:vision-video-restoration-evaluation}

多帧方法的收益和成本都取决于输入范围。若一个方法看未来十帧，另一个必须立即输出，单帧质量差异就同时包含信息条件差异。合理的比较先统一序列、退化和用途，再说明算法怎样使用时间。

\subsection{比较协议与成本的统一}
\label{sec:vision-video-restoration-protocol}

超分需要固定空间倍率、下采样与模糊方式、裁边和色彩空间；插帧需要固定两端间隔、目标时刻以及是否含先前生成帧；压缩需要固定输入分辨率、帧率、关键帧间隔和码流统计。不同源视频划分与近重复片段也会影响结果，不能把相邻帧泄漏到不同数据划分。

对帧宽$w$、高$h$、共$n$帧、时长$T$秒的视频，完整码流为$B$字节时，可分别报告
\[
  \text{平均比特/像素}=\frac{8B}{nwh},
  \qquad
  \text{平均比特率}=\frac{8B}{T}\ \mathrm{bit/s}\text{。}
\]
例如10秒、30帧每秒、$640\times360$的视频，完整码流恰为300000字节，则平均码率为240000 bit/s，平均约0.03472 bit/像素。这只是构造的成本计算，不对应任何模型实测表现。

计算延迟、等待后帧的延迟和显示缓冲分别测量。吞吐量每秒处理多少帧，与某帧从拍摄到输出等待多久是不同量；批处理吞吐很高，也可能让首帧等待很久。报告时还需给出硬件、精度、批量、输入长度和峰值内存，避免只比较网络某个局部模块的耗时。

\subsection{内容、时间连续性与用途的核验}
\label{sec:vision-video-temporal-quality}

PSNR、SSIM及感知距离继续评价单帧，但视频还要检查闪烁、重影、拖尾和字符随时间变化。相邻帧正常运动时像素本来就不同，直接惩罚$\widehat{\bm{X}}_t-\widehat{\bm{X}}_{t-1}$会鼓励画面停滞。时间比较应补偿真实运动，并排除遮挡、新显露和越界区域。

一种诊断量是在可信可见区域中比较当前输出与邻帧输出的对齐差。若$M_t(\bm{p})$为有效性掩码，$\mathcal{W}$使用独立估计或标注的对应，则可计算
\begin{equation}
  E_{\mathrm{temp}}=
  \frac{\sum_{t,\bm{p}}M_t(\bm{p})
    \left\lVert\widehat{\bm{X}}_t(\bm{p})
      -\mathcal{W}(\widehat{\bm{X}}_{t-1},
        \bm{u}_{t\to t-1})(\bm{p})\right\rVert_1}
       {c\sum_{t,\bm{p}}M_t(\bm{p})}\text{，}
  \label{eq:vision-video-temporal-error}
\end{equation}
其中$c$为通道数，分母非零时才有定义。这是归一化的诊断形式，并非覆盖全部视频质量的标准；光流不准、曝光变化或掩码漏判都可能影响它。计算必须说明对应和有效性来自哪里。

时间差小也不保证内容正确。每帧都输出同一张模糊灰图，可以很稳定却丢失全部细节；每帧都把车牌同一字符改错，也可能保持时间一致。因此，应同时检查与参考内容的关系及实际用途。逐帧阅读车牌、回查原观测、比较置信度和字符串变化，能识别“单帧好看、播放时字符反复变换”的失败。

有配对清晰视频时，还可比较输出与真值的运动补偿差异；没有真值时，主观连续性、下游任务和原始证据检查更重要，但不能将这些替代检查描述为已证明像素恢复正确。把失败追溯到取样、可靠性、状态传播、生成先验或编码参考，才有助于改进处理方案。

\subsection{质量、码率与延迟的选择}
\label{sec:vision-video-quality-rate-latency}

固定源序列后，用多个配置形成质量与码率、质量与等待时间的对应关系。若一个配置更省码率却必须等待整段视频，另一个略差但能即时输出，选择取决于用途。观看容忍适量细节损失，车牌读取更在意字符，测量还要求原始观测和时间坐标可追溯。

表\ref{tab:vision-video-tradeoff}汇总不同处理职责最容易遗漏的条件。它不是方法排名，而是解释为什么同一“清晰视频”目标需要多项核验。真实时变退化、大运动、遮挡、切镜和极低码率都可能改变取舍；更多帧或更复杂网络只有在保留所需信息、满足等待约束时才有实际价值。
\begin{table}[htbp]
  \centering
  \begin{tabularx}{\linewidth}{@{}lXX@{}}
    \toprule
    处理职责 & 主要新增信息 & 必须同时检查\\
    \midrule
    多帧复原 & 邻帧互补内容 & 对应、可见性与错误传播\\
    空间超分 & 位移形成的采样差异 & 真实细节与先验生成细节\\
    时间插帧 & 两端的时间约束 & 后帧等待、遮挡与运动歧义\\
    帧间压缩 & 已解码参考的可预测内容 & 码流、参考一致与随机访问\\
    \bottomrule
  \end{tabularx}
  \caption{视频处理的不同收益依赖不同信息条件。比较应固定输入与用途，并把右列条件和单帧质量一起检查，不能据某一项改善推断整个系统已满足要求。}
  \label{tab:vision-video-tradeoff}
\end{table}
\FloatBarrier

\section{本章小结}
\label{sec:vision-video-restoration-summary}

邻帧只有在对应正确、内容可用时才能补充当前画面。反向取样解决到哪里读取，可靠性和注意力决定读来的信息如何参与，状态传播则决定哪些历史仍能影响后续输出。空间超分与时间插帧扩展了不同方向的采样，也分别引入未直接观测内容的估计。

编码进一步要求共享信息能在解码端重现。运动和残差的比特要共同分配，预测只能依赖已解码参考，完整码流须能独立恢复。明确这些依赖之后，才可以在内容忠实、播放连续、码率和等待之间作出有依据的选择；画面更锐利、更平滑或张量更小，都不足以单独证明任务完成。
